% preamble/part.title.tex
% Part 扉页设计（蓝金配色，分析学笔记）
%
% 与章节扉页保持一致的蓝金风格，布局为章节页的"镜像"：
%   - 顶部全宽深蓝横条（与章节页中部的彩条完全一致）
%   - 左上：大号金色罗马数字 + Part 标题
%   - 右上：简介区域
%   - 金色装饰分割线
%   - 下方：整页宽图片填满剩余空间
%
% 接口：
%   \setpartimage{<path>}   -- 设置下方大图
%   \setpartintro{<text>}   -- 设置右上简介文字
%
% 用法示例：
%   \setpartimage{figure/part1.png}
%   \setpartintro{这一部分从 Peano 公理出发......}
%   \part{单变量分析学}
% ============================================================

\RequirePackage{tikz}
\RequirePackage{graphicx}
\RequirePackage{etoolbox}
\RequirePackage{setspace}

\makeatletter

% -------- 每 Part 图片接口 --------
\gdef\partimage{}
\newcommand{\setpartimage}[1]{\gdef\partimage{#1}}

% -------- Part 简介接口 --------
\gdef\partintro{}
\newcommand{\setpartintro}[1]{\gdef\partintro{#1}}

% -------- 版式参数 --------
% 顶部蓝色横条（与章节页中部彩条一致）
\newlength{\PartRibbonHeight}
\setlength{\PartRibbonHeight}{2.35cm}

% 横条下方到内容区域的间距
\newlength{\PartContentTop}
\setlength{\PartContentTop}{\dimexpr\PartRibbonHeight+1.2cm\relax}

% 金色分割线的垂直位置
\newlength{\PartDividerY}
\setlength{\PartDividerY}{\dimexpr\PartRibbonHeight+6.8cm\relax}

% -------- 左上 Part 标题 block --------
\newcommand{\PartTitleBlock}[2]{%
    \ifdefempty{#1}{%
        % 无编号：菱形装饰
        \fill[GoldMain!55]
            ([xshift=1.8cm,yshift=-\PartContentTop]current page.north west)
            -- ++(0.22,0.22) -- ++(-0.22,0.22) -- ++(-0.22,-0.22) -- cycle;
    }{%
        % Part 罗马数字（大号金色）
        \node[anchor=north west,text=GoldMain!72,font=\fontsize{52}{62}\bfseries]
            at ([xshift=1.8cm,yshift=-\PartContentTop]current page.north west)
            {#1};
        % "Part" 标注
        \node[anchor=north west,text=GoldMain!48,font=\Large\bfseries]
            at ([xshift=1.8cm,yshift=-\PartContentTop-3.1cm]current page.north west)
            {Part};
    }%
    % Part 标题（深蓝色）
    \node[anchor=north west,text=BlueDark,font=\fontsize{22}{28}\bfseries,
          text width=8cm,align=left]
        at ([xshift=1.8cm,yshift=-\PartContentTop-4.4cm]current page.north west)
        {#2};
}

% -------- 右上简介 block --------
\newcommand{\PartIntroBlock}{%
    \ifdefempty{\partintro}{}{%
        % "简介" 标题
        \node[anchor=north east,text=BlueMain!80,font=\Large\bfseries]
            at ([xshift=-1.8cm,yshift=-\PartContentTop]current page.north east)
            {简\quad 介};
        % 简介正文
        \node[anchor=north east,text=BlueDark!55,font=\large\itshape,
              text width=7.4cm,align=right,inner sep=0pt]
            at ([xshift=-1.8cm,yshift=-\PartContentTop-1.5cm]current page.north east)
            {\setstretch{1.45}\partintro};
    }%
}

% -------- Part 扉页 --------
%   #1 = Part 编号（罗马数字，如 I, II），空串表示无编号
%   #2 = Part 标题
\newcommand{\PartTitlePage}[2]{%
    \thispagestyle{empty}%
    \vspace*{-\baselineskip}%
    % 将编号归一化为全局宏 \partnum：有编号时为 \thepart 的内容，
    % 无编号时为空宏。这样 \PartTitleBlock 里的 \ifdefempty 永远
    % 收到的是合法控制序列（而非空分组 {}），避免 \part* 渲染报错。
    \gdef\partnum{#1}%
    \begin{tikzpicture}[remember picture,overlay]
        % ============================================================
        %  第一层：顶部全宽深蓝横条（与章节页中部彩条完全一致）
        % ============================================================
        \fill[BlueMain!94]
            (current page.north west) rectangle
            ([yshift=-\PartRibbonHeight]current page.north east);
        % 左侧金色装饰竖条
        \fill[GoldMain!75]
            (current page.north west) rectangle
            ([xshift=0.18cm,yshift=-\PartRibbonHeight]current page.north west);
        % 右侧金色装饰竖条
        \fill[GoldMain!75]
            ([xshift=-0.18cm]current page.north east) rectangle
            ([yshift=-\PartRibbonHeight]current page.north east);

        % ============================================================
        %  第二层：左上 —— Part 编号 + 标题
        % ============================================================
        \PartTitleBlock{\partnum}{#2}

        % ============================================================
        %  第三层：右上 —— 简介区域
        % ============================================================
        \PartIntroBlock

        % ============================================================
        %  第四层：左右区域之间的金色装饰虚线
        % ============================================================
        \draw[GoldMain!20,line width=0.5pt,dash pattern=on 3pt off 5pt]
            ([xshift=11.0cm,yshift=-\PartContentTop]current page.north west) --
            ([xshift=11.0cm,yshift=-\PartDividerY+0.3cm]current page.north west);

        % ============================================================
        %  第五层：金色装饰分割线（两端圆点 + 中央菱形）
        % ============================================================
        \draw[GoldMain!42,line width=0.7pt]
            ([yshift=-\PartDividerY]current page.north west) --
            ([yshift=-\PartDividerY]current page.north east);
        \fill[GoldMain!50]
            ([yshift=-\PartDividerY]current page.north west) circle (1.8pt);
        \fill[GoldMain!50]
            ([yshift=-\PartDividerY]current page.north east) circle (1.8pt);
        \fill[GoldMain!60]
            ([yshift=-\PartDividerY]current page.north)
            -- ++(0.13,0.13) -- ++(-0.13,0.13) -- ++(-0.13,-0.13) -- cycle;

        % ============================================================
        %  第六层：下方全宽图片（填满剩余空间）
        % ============================================================
        \ifdefempty{\partimage}{%
            % 无图片时显示浅蓝占位渐变
            \shade[top color=BlueMain!12,bottom color=BlueMain!3]
                ([yshift=-\PartDividerY]current page.north west) rectangle
                (current page.south east);
            \node[align=center,text=BlueDark!40,font=\large\itshape]
                at ([yshift=-5.5cm]current page.center)
                {请在 \texttt{\textbackslash part} 前用
                 \texttt{\textbackslash setpartimage} 设置图片};
        }{%
            \node[anchor=north west,inner sep=0pt,outer sep=0pt]
                at ([yshift=-\PartDividerY]current page.north west)
                {\includegraphics[
                    width=\paperwidth,
                    height=\dimexpr\paperheight-\PartDividerY\relax
                ]{\partimage}};
        }%
    \end{tikzpicture}%
    \clearpage%
    % 双面模式下的空白页
    \if@twoside
        \if@openright
            \null
            \thispagestyle{empty}%
            \newpage
        \fi
    \fi
}

% ============================================================
%  重定义 \@part（编号 Part）
% ============================================================
\def\@part[#1]#2{%
    \ifnum \c@secnumdepth >-2\relax
        \refstepcounter{part}%
        \addcontentsline{toc}{part}{\thepart\hspace{1em}#1}%
    \else
        \addcontentsline{toc}{part}{#1}%
    \fi
    \markboth{}{}%
    \PartTitlePage{\thepart}{#2}%
    \@endpart
}

% ============================================================
%  重定义 \@spart（无编号 Part，即 \part*）
%  —— 与 \part{} 完全一致：渲染整页扉页、写入目录；
%     区别仅在于不显示罗马数字（\PartTitleBlock 收到空编号，
%     自动以菱形装饰代替）。
%  —— 说明：\secdef 对星号形式直接展开为 \@spart{标题}，
%     不经过 \@dblarg，故 \@spart 只收单个必选参数，
%     不支持 \part*[目录文字]{标题} 这种可选参数写法。
%  —— 原模板的 \@spart 漏写了 \addcontentsline，导致 \part*
%     不进目录；这里补上（不带编号），使行为与 \part{} 对齐。
% 用法：
%   \setpartimage{figure/partX.png}
%   \setpartintro{……}
%   \part*{一些短篇——分析篇}        % 目录与扉页均无编号
% ============================================================
\def\@spart#1{%
    \addcontentsline{toc}{part}{#1}%
    \markboth{}{}%
    \PartTitlePage{}{#1}%
    \@endpart
}

\def\@endpart{%
    \if@tempswa
        \twocolumn
    \fi
}

\makeatother
